\chapter*{}
\begin{center}
    {
    \Large{
        \textbf{SISTEMA DE MONITORAMENTO ESTRUTURAL E HÍDRICO DE BARRAGENS UTILIZANDO IOT E VISÃO COMPUTACIONAL EM EDGE COMPUTING}
        }
    }
\end{center}

\vspace{1cm}

\begin{flushright}
    Autor: Gladerson Pereira Belizio da Silva\\
\end{flushright}

\vspace{1cm}

\begin{center}
    \Large{\textsc{\textbf{Resumo}}}
\end{center}

\noindent 
A gestão dos recursos hídricos e a garantia da segurança de barragens são aspectos críticos para a sustentabilidade do semiárido brasileiro. Nesse contexto, este trabalho propõe uma arquitetura custo-efetiva, fundamentada em Internet das Coisas (IoT) e Inteligência Artificial (IA), aplicada ao monitoramento da Barragem de Oiticica (RN). O sistema integra sensores de nível baseados em tecnologia radar e processamento de imagens na borda (Edge Computing), empregando microcomputadores de placa única e técnicas de visão computacional baseadas em aprendizado profundo para a detecção autônoma de anomalias em estruturas de concreto. A visualização dos dados é consolidada em uma plataforma de supervisão que correlaciona a telemetria em tempo real a uma representação tridimensional georreferenciada da infraestrutura monitorada.

\noindent\textit{Palavras-chave}: IoT; Visão Computacional; Edge Computing; Monitoramento de Barragens; Aprendizado Profundo.
